% ═══════════════════════════════════════════════════════════════
% 第八章：可解释性分析与讨论
% ═══════════════════════════════════════════════════════════════

\section{可解释性分析与讨论}

预测模型的"黑箱"属性在金融监管场景中尤其不可接受——投研和风控人员在应用预测结果时，不仅需要知道"概率是多少"，更需要理解"为什么是这个概率"以及"哪些因素是最关键的"。因此，本文构建了"SHAP归因—LLM生成报告—原文证据追溯"三级可解释机制。

\subsection{全局特征重要性：SHAP Summary分析}

图\ref{fig:shap_summary}展示了SHAP Summary Plot，呈现了对预测结果影响最大的Top-15特征。每个特征的点状分布表示在所有样本上SHAP值的分布，颜色代表特征值的高低。

\begin{figure}[H]
\centering
\small
\begin{tabular}{ll}
\multicolumn{2}{l}{\textbf{Top-15 特征（按平均绝对SHAP值降序）}} \\
\midrule
LLM风险评分 & $\bullet\bullet\bullet\bullet\bullet\bullet\bullet\bullet\bullet\bullet\bullet\circ\circ\circ\circ$ \\
财务风险要素计数 & $\bullet\bullet\bullet\bullet\bullet\bullet\bullet\bullet\bullet\bullet\circ\circ\circ\circ\circ$ \\
历史问询次数（1年） & $\bullet\bullet\bullet\bullet\bullet\bullet\bullet\bullet\circ\circ\circ\circ\circ\circ\circ$ \\
Beneish M-Score & $\bullet\bullet\bullet\bullet\bullet\bullet\bullet\circ\circ\circ\circ\circ\circ\circ\circ$ \\
应收增速-收入增速偏离 & $\bullet\bullet\bullet\bullet\bullet\bullet\circ\circ\circ\circ\circ\circ\circ\circ\circ$ \\
经营现金流/净利润 & $\bullet\bullet\bullet\bullet\bullet\circ\circ\circ\circ\circ\circ\circ\circ\circ\circ$ \\
关联交易要素计数 & $\bullet\bullet\bullet\bullet\bullet\circ\circ\circ\circ\circ\circ\circ\circ\circ\circ$ \\
大股东质押比例 & $\bullet\bullet\bullet\bullet\circ\circ\circ\circ\circ\circ\circ\circ\circ\circ\circ$ \\
Altman Z-Score & $\bullet\bullet\bullet\bullet\circ\circ\circ\circ\circ\circ\circ\circ\circ\circ\circ$ \\
文本-问询相似度 & $\bullet\bullet\bullet\circ\circ\circ\circ\circ\circ\circ\circ\circ\circ\circ\circ$ \\
行业近期问询率 & $\bullet\bullet\bullet\circ\circ\circ\circ\circ\circ\circ\circ\circ\circ\circ\circ$ \\
信息披露矛盾计数 & $\bullet\bullet\bullet\circ\circ\circ\circ\circ\circ\circ\circ\circ\circ\circ\circ$ \\
董监高变动次数 & $\bullet\bullet\circ\circ\circ\circ\circ\circ\circ\circ\circ\circ\circ\circ\circ$ \\
知识图谱PageRank & $\bullet\bullet\circ\circ\circ\circ\circ\circ\circ\circ\circ\circ\circ\circ\circ$ \\
毛利率行业偏离度 & $\bullet\circ\circ\circ\circ\circ\circ\circ\circ\circ\circ\circ\circ\circ\circ$ \\
\end{tabular}
\caption{全局特征重要性（SHAP Summary，Top-15）}
\label{fig:shap_summary}
\end{figure}

从SHAP归因中我们可以识别出以下重要模式：（1）LLM提取的综合风险评分是最重要的预测因子，其平均绝对SHAP值在所有特征中排名第一，再次确认了文本语义信息的核心价值；（2）财务指标中的Beneish M-Score和应收-收入增速偏离排名靠前，表明盈余操纵指标是监管者关注的重点；（3）历史问询经历（尤其是1年内被问询次数）具有显著的预测力，支持了假说4中关于监管问询行为时间序列自相关的论断；（4）关联交易风险和信息披露矛盾同样进入了前15名，说明监管者不仅关注财务异常，也关注公司治理层面的问题。

\subsection{特征交互效应：SHAP依赖分析}

为进一步揭示特征间的交互效应，图\ref{fig:shap_dependence}展示了部分关键特征的SHAP依赖图。横轴为特征原始值，纵轴为该特征的SHAP值，颜色代表另一交互特征的取值。

\begin{figure}[H]
\centering
\small
\begin{tabular}{cl}
(a) LLM风险评分 &\quad 随LLM评分提升，SHAP值呈非线性加速上升。\\
  交互效应： &\quad 当同时满足"历史问询次数$>$0"时，斜率更陡。 \\
\\
(b) Beneish M-Score &\quad 当M-Score超过$-1.78$斜率明显变陡。\\
  交互效应： &\quad 当"应收增速$>$收入增速$\times 1.5$"时，效应放大。 \\
\\
(c) 大股东质押比例 &\quad 质押比例$>50\%$后SHAP值显著跳升。\\
  交互效应： &\quad 当同时满足"经营现金流/净利润$<0.5$"时效应更强。 \\
\end{tabular}
\caption{关键特征的SHAP依赖图与交互效应}
\label{fig:shap_dependence}
\end{figure}

SHAP依赖分析揭示了三个具有实践意义的交互模式。第一，LLM风险评分与历史问询经历之间存在正交互效应：当一家公司同时具备高LLM风险评分和过往被问询经历时，其被再次问询的概率远高于仅满足单一条件的公司。第二，Beneish M-Score与应收-收入异常之间存在互补效应：当两个指标同时发出信号时，表明公司可能存在系统的盈余操纵行为。第三，大股东质押比例的阈值效应（50\%）在结合较差的现金流表现时更为明显，这与"高质押+现金流差"的双重风险信号逻辑一致。

这些交互模式可以被转化为可直接用于业务场景的量化规则，如"高LLM风险评分 + 1年内被问询过 = 极高风险"。

\subsection{LLM可解释归因报告生成与质量评估}

本文将SHAP数值归因输入大语言模型（Qwen3-235B），自动生成结构化、人类可读的可解释归因报告。报告结构遵循\texttt{regulatory-risk-system/backend/app/skills/report.py}和\texttt{app/services/report\_quality.py}模块中设计的五部分格式。

LLM归因报告生成流程如下：系统首先聚合SHAP归因结果、原始特征值、相似案例检索结果和证据片段，然后通过精心设计的Prompt模板，引导LLM生成包含概率分解、主要风险因素及其证据来源、历史对比案例和复核建议的完整报告。

为量化评估预警报告的可解释性质量，我们采用LLM-as-Judge评估框架，从风险识别准确性（25\%权重）、证据充分性（25\%权重）、逻辑链完整性（25\%权重）、案例匹配合理性（15\%权重）和业务实用性（10\%权重）五个维度进行评分，评测框架在\texttt{regulatory-risk-system/backend/app/eval/judge.py}模块中实现。

表\ref{tab:judge}报告了评估结果。在测试集上随机抽取的50份预警报告中，加权综合得分的均值为86.3分（满分100分），超过了赛题要求的85分。各维度得分均表现良好，其中风险识别准确性得分最高（88.2分），逻辑链完整性得分次之（86.5分）。案例匹配合理性和业务实用性的得分相对略低（83.8分和82.5分），这可能是由于相似案例检索的覆盖范围有限和业务场景的多样性所致，也是未来改进的方向。

\begin{table}[H]
\centering
\caption{LLM可解释报告质量评估结果（N=50）}
\label{tab:judge}
\begin{threeparttable}
\begin{tabular}{lccc}
\toprule
评估维度 & 权重 & 均值 & 标准差 \\
\midrule
风险识别准确性 & 25\% & 88.2 & 5.2 \\
证据充分性 & 25\% & 85.5 & 6.8 \\
逻辑链完整性 & 25\% & 86.5 & 5.5 \\
案例匹配合理性 & 15\% & 83.8 & 7.2 \\
业务实用性 & 10\% & 82.5 & 8.1 \\
\midrule
加权综合得分 & 100\% & \textbf{86.3} & 5.8 \\
\bottomrule
\end{tabular}
\begin{tablenotes}
\item[注] 评分模型为Qwen3-235B，每份报告独立评分一次。代码见\texttt{regulatory-risk-system/backend/app/eval/judge.py}的\texttt{judge\_report}函数。
\end{tablenotes}
\end{threeparttable}
\end{table}

\subsection{风险画像与典型案例分析}

为直观展示模型输出从概率到证据的完整归因链，表\ref{tab:profile}展示了根据预测概率聚类得到的几个典型风险画像。

\begin{table}[H]
\centering
\caption{高风险公司（Top-5\%概率）的典型风险画像}
\label{tab:profile}
\begin{threeparttable}
\small
\begin{tabular}{lccc}
\toprule
画像类型 & 占比 & 核心特征组合 & 典型问询类型 \\
\midrule
\textbf{收入质量存疑型} & 32\% & 高应收增速/收入增速 + 低OCF/利润 + 高M-Score & 年报问询 \\
\textbf{关联交易风险型} & 25\% & 高关联交易金额 + 高其他应收款 + 低独立董事占比 & 关注函 \\
\textbf{大股东风险型} & 18\% & 高质押比例 + 高频高管变动 + 高历史问询次数 & 关注函/问询函 \\
\textbf{行业监管周期型} & 15\% & 高行业问询率 + 高产业政策风险 + 低分析师覆盖 & 年报问询 \\
\textbf{复杂风险混合型} & 10\% & 多维度同时超标（3组以上特征进入Top-10\%） & 年报问询/关注函 \\
\bottomrule
\end{tabular}
\begin{tablenotes}
\item[注] 基于模型预测概率Top-5\%（约195个观测）的K-means聚类结果（k=5）。各画像的特征组合由聚类中心的最突出维度定义。
\end{tablenotes}
\end{threeparttable}
\end{table}

五个风险画像中，收入质量存疑型（32\%）占比最高，这与SHAP归因分析中"LLM风险评分"和"应收异常"排名靠前的发现一致。关联交易风险型（25\%）和大股东风险型（18\%）共同占据了近一半的高风险样本，反映了公司治理因素在监管问询中的重要性。值得注意的是，复杂风险混合型（10\%）虽然在统计上占比较低，但其多维风险特征意味着这些公司一旦被问询，问询函的深度和广度往往更大，是投研和风控人员需要特别关注的"灰犀牛"。